\subsection{基本态射}\label{sec:17d}
令 $M$ 为带角正合辛流形，$V$ 为其中带标架的正合 Lagrangian 球面。为方便利用\stref{笔记}{17.8}中上链层次的消失结果，假设 $2c_1(M,V)$ 消失。置 $L_0=V$，另取两个正合 Lagrangian 子流形 $L_1,L_2$，要求 $L_2$ 与 $\tau_V(L_1)$ 同痕；更准确地，还要指定一个同痕选择。为子流形对 $(L_0,L_1)$、$(L_0,L_2)$、$(L_1,L_2)$ 固定 Floer 数据，并固定定义乘积 $\mu^2:CF^{\mathrm{pr}}(L_1,L_2)\otimes CF^{\mathrm{pr}}(L_0,L_1)\to CF^{\mathrm{pr}}(L_0,L_2)$ 所需的扰动数据。为便于与代数理论比较，用 $\Y_0,\Y_1,\Y_2$ 表示 $\F^{\mathrm{pr}}(M)$ 中由 $L_0,L_1,L_2$ 给出的对象。

固定圆盘上的 Lefschetz 纤维化 $\pi:E\to S$，纤维为 $M$，只有一个临界点，消失圈为 $V$。更准确地，要求该纤维化带有标点 $*\in\partial S$，并将该点上方纤维 $E_*$ 与 $M$ 认同；还带有从 $*$ 出发的消失路径 $\gamma$，使 $V_\gamma=V$。为简化技术，进一步假设纤维化在 $*$ 附近局部平凡。$E$ 将是后续全部构造的基本组成部分，这些构造基本上是往它上面粘合各种平凡部分。

\StatementNumber{9}\begin{note}\label{st:17.9}
回顾，对任意给定的 $M,V$，都存在这样的纤维化，而且可以“手工”构造；概要见\secref{16e}，完整论述见\bcite{131}引理 1.10。仔细关注构造细节，还能获得更多控制：辛联络在临界点的小邻域之外平坦；任意消失路径的消失圈均等于 $V$；任意绕临界值一周的回路的单值变换，都是一个 Dehn 扭转 $\tau_V$，而不只是与之同痕。不过，这些性质对于我们的构造都并非必要。
\end{note}

作为第一个也是最初等的应用，定义 Floer 上闭链
\begin{equation}\label{eq:17.12}
\boxed{\begin{gathered}
c\in CF^{\mathrm{pr}}(L_1,L_2),\\ \mu^1(c)=0.
\end{gathered}}
\end{equation}
\SourcePage{252}
为此，取给定的 $E\to S$ 并去掉 $*$。所得 Lefschetz 纤维化 $E^c\to S^c$ 有一个无穷远点，将其视为入向点，记为 $\zeta^{c,-}$，并选择条带状端 $\epsilon_{\zeta^{c,-}}$。存在逐纤维同痕意义下典范的移动 Lagrangian 边界条件 $F^c$，定义如下：在 $\epsilon_{\zeta^{c,-}}(s,0)$，$s\ll0$，上方的纤维中从 $L_1$ 出发，通过平行移动沿 $\partial S^c$ 几乎绕行一周。结果与 $\tau_V(L_1)$ 同痕，因而也与 $L_2$ 同痕；用该同痕将其连续延伸，使在 $\epsilon_{\zeta^{c,-}}(s,1)$，$s\ll0$，上方的纤维中等于 $L_2$。

\sr{fig:17.1}{图~17.1}示意这一构造。图形记号的含义是：去掉奇异纤维，再沿虚线切开底空间，得到平凡辛纤维化。随后将单值变换看成切口两侧之间的认同。在此类图中，通常不区分同痕的对象，所以直接将单值变换写为 $\tau_V$。
\SourceFigure{figure-17-1}{0.37176883}{17.1}
于是，\eref{17.12}就是相应不变量 $C\Phi_{E^c/S^c}$，取某个与已有 Floer 数据选择相容的正则相对扰动数据。

第二项数据是 $\mu^2(c,\cdot)$ 与零映射之间的链同伦，即
\begin{equation}\label{eq:17.13}
\boxed{\begin{gathered}
k:CF^{\mathrm{pr}}(L_0,L_1)\longrightarrow CF^{\mathrm{pr}}(L_0,L_2),\\
\mu^1(k(\cdot))+k(\mu^1(\cdot))+\mu^2(c,\cdot)=0.
\end{gathered}}
\end{equation}
如所预期，这是由依赖于参数 $r\in[0,1]$ 的带条带状端 Lefschetz 纤维化族构造的。在一端 $r=0$，取三标点圆盘上的平凡纤维化，带 Lagrangian 边界条件 $(L_0,L_1,L_2)$。选取适当扰动数据，其关联不变量正好给回 $\mu^2$。现取两个出向端中标记为 $(L_1,L_2)$ 的那个，与 $E^c\to S^c$ 的入向端粘合。粘合长度应充分大，使零维伪全纯截面模空间的粘合公式在上链层次成立。所得纤维化记为 $E^{k,0}\to S^{k,0}$，配备移动边界条件 $F^{k,0}$
\SourcePage{253}
和相对扰动数据，其相应不变量恰为 $\mu^2(c,\cdot)$。\sr{fig:17.2}{图~17.2}给出两个图景，只在我们设想中切开底空间的方式上有所不同。在另一端 $r=1$，按如下方式进行。
\SourceFigure{figure-17-2}{0.48184563}{17.2}

取 $Z$ 上的平凡纤维化。赋予移动 Lagrangian 边界条件，在 $\R\times\{0\}$ 上等于 $L_0$，在 $\R\times\{1\}$ 上实现从 $L_2$ 到 $\tau_V(L_1)$ 的同痕。这意味着，两端关联的 Lagrangian 子流形对分别是 $(L_0,L_2)$ 与 $(L_0,\tau_V(L_1))$。对于后一对，取 $(L_0,L_1)$ 的 Floer 数据在 $\tau_V$ 下的像，同时记住 $\tau_V(L_0)=\tau_V(V)=V=L_0$。还要标记边界点 $z^+\in\R\times\{0\}$。现取原来的 $\pi:E\to S$，赋予标准边界条件 $F$，使其在 $*$ 上方纤维中等于 $L_0$，且在该点附近局部恒定。最后，沿 $z^+$ 与 $z^-=*$ 取边界连通和，将两个纤维化粘合，结果记为 $E^{k,1}\to S^{k,1}$，见\sr{fig:17.3}{图~17.3}。这是\stref{引理}{17.7}讨论的过程的一个实例，事实上\stref{笔记}{17.8}也适用；因此，选择适当
\SourcePage{254}
小的粘合参数，可保证 $E^{k,1}$ 的所有零维伪全纯截面模空间，在 $n>1$ 时为空，在 $n=1$ 时包含偶数个点。
\SourceFigure{figure-17-3}{0.3489227}{17.3}

直接比较\sr{fig:17.2}{图~17.2}和\sr{fig:17.3}{图~17.3}可见，确有单参数纤维化族 $E^{k,r}\to S^{k,r}$ 及边界条件 $F^{k,r}$，连接我们在 $r\in\{0,1\}$ 时构造的对象。这些纤维化的无穷远结构，即条带状端上的结构，对所有 $r$ 相同。选择一族相对扰动数据，在 $r=0,1$ 时与之前构造的相同，并考虑所得带参数伪全纯截面模空间。计数其中孤立点，就定义同态 $k:CF^{\mathrm{pr}}(L_0,L_1)\cong CF^{\mathrm{pr}}(L_0,\tau_V(L_1))\to CF^{\mathrm{pr}}(L_0,L_2)$，它是 $C\Phi_{E^{k,0}/S^{k,0}}$ 与 $C\Phi_{E^{k,1}/S^{k,1}}$ 之间的链同伦，故恰具有所需性质\eref{17.13}。

构造 $c$ 时必须作一些选择，因为边界条件及扰动数据都不完全典范。即使指定了它们，定义族 $E^{k,r}\to S^{k,r}$ 时仍有进一步选择，又会影响 $k$。但两种情况下，参数空间都道路连通，故能在任意两种可能之间光滑移动。具体地，若 $(c,k)$ 与 $(\tilde c,\tilde k)$ 用两种不同选择定义，则可考虑多一个参数的模空间，从而找到
\begin{equation}\label{eq:17.14}
\boxed{\begin{gathered}
\gamma\in CF^{\mathrm{pr}}(L_1,L_2),\\
\kappa:CF^{\mathrm{pr}}(L_0,L_1)\longrightarrow CF^{\mathrm{pr}}(L_0,L_2),\\
\mu^1(\gamma)+c+\tilde c=0,\\
\mu^1(\kappa(\cdot))+\kappa(\mu^1(\cdot))+\mu^2(\gamma,\cdot)+k+\tilde k=0.
\end{gathered}}
\end{equation}
这只是此前考察的直接扩展，因此不
\SourcePage{255}
解释细节。下面改用第\sr{sec:5}{5 节}的代数理论来解释所得结果。

\StatementNumber{10}\begin{note}\label{st:17.10}
注意，Floer 理论框架尚未引入分次，而且我们在 $\operatorname{char}(\K)=2$ 下工作，符号是平凡的；上面的公式清楚地体现了二者的缺席。不过，\stref{笔记}{9.7}在此也适用：首先，为不分次 $\Ai$ 范畴改写代数理论相当容易；其次，当前背景只是暂时的，后面会加以修改，纳入适当的分次与符号。因此，我们可以直接使用第\sr{sec:5}{5 节}的结果。
\end{note}

具体地，\stref{引理}{5.3}(i) 结合\eref{17.12}与\eref{17.13}表明，$(c,k)$ 在 $\Ai$ 模的上同调范畴中给出态射 $T_{\Y_0}(\Y_1)\to\Y_2$。类似地，\stref{引理}{5.3}(ii) 与\eref{17.14}说明该态射典范；即它与用来定义 $c,k$ 的 Lefschetz 纤维化和模空间的定义中的辅助选择无关。\znote{247}{这一段是本节的落点：$c$ 是上闭链、$k$ 把 $\mu^2(c,\cdot)$ 同伦到零，两者合起来正是“$\mathcal T_{Y_0}(Y_1)\to Y_2$ 是 $\Ai$ 模态射”在链级的完整内容，用第 5 节的语言即给出同态 $(\Id,c,k,\ldots)$。之所以要求 $L_2$ 与 $\tau_V(L_1)$ 同痕，是因为 $c$ 的边界条件沿底空间绕一圈后正好把 $L_1$ 送到 $\tau_V(L_1)$；若 $L_2$ 不与之同痕，$c$ 便没有落点。这一构造是 §(17g)–(17j) 反复使用的部件，也是定理 17.20 的出发点；该定理断言：在拟等价意义下，$\Tw\F^\to(V^\#)$ 是 $V$ 的 Hurwitz 等价类的不变量，其证明的代数核心正是“几何 Dehn 扭转 $\tau_V$ 对应扭转函子 $T_{Y_0}$”。[来源：原书 §(17d) 与 §(17g)，印刷页 244–247；定理编号见原书目录。]}
